\section{1.1 Motivation}\label{motivation}

A language model with no persistent memory cannot reliably carry
preferences, commitments, or corrections across sessions. A language
model with persistent memory is risky for the opposite reason: stale
facts can keep influencing behavior after the user has corrected them,
contradictory records can accumulate silently, and a system can act on
information it should have refused to write in the first place. The
engineering objective is therefore not ``remember as much as possible''
but to govern the memory write as a decision under explicit policy.

This observation changes where control must live. Prompting a model to
``be careful'' is insufficient because memory writes have side effects
that outlive the prompt. A new candidate fact must be classified against
the existing store --- is it independent, a duplicate, or a
contradiction of a prior record? --- and the resulting action must be
explicit and auditable: keep the new record alongside the old, supersede
the old record while preserving its lineage, or abstain because the
evidence is too ambiguous to commit either way. When a system cannot
express the third option, every conflict becomes a silent overwrite or a
silent duplicate.

Existing evaluations mostly treat memory as an answer-retrieval problem:
given a long interaction history, can the model answer a later question
\citep{wu2025longmemeval, maharana2024locomo}? That view is valuable but
incomplete: end-to-end question accuracy cannot attribute a failure to
the write decision that caused it, and it cannot see governance failures
that happen not to affect the sampled questions. A governance evaluation
must score the decision boundary itself.

\section{1.2 Research problem}\label{research-problem}

\begin{quote}
How can a language-agent memory system make and audit policy-sensitive
write decisions --- keep, supersede, abstain --- and how can those
decisions be evaluated without conflating governance correctness with
retrieval or answer quality?
\end{quote}

The problem decomposes into five concerns: \textbf{representation}
(typed memory records with provenance and lineage), \textbf{decision}
(an interpretable, abstention-aware action policy), \textbf{execution}
(supersession must actually mark the prior record inactive and preserve
its lineage), \textbf{measurement} (decision-level metrics with correct
denominators, plus state metrics such as lineage integrity and stale
retrieval, distinct from answer metrics), and \textbf{honest reporting}
(every number carries a claim-eligibility label; contaminated or patched
runs are labeled as such).

\section{1.3 Research questions}\label{research-questions}

\textbf{RQ1 --- Architecture.} What controller decomposition makes
memory-write decisions explicit, interpretable, and auditable, and
prevents a language model from mutating memory outside a governed action
vocabulary?

\textbf{RQ2 --- Diagnostic.} Can a compact, author-designed diagnostic
slice expose memory-governance failures --- invalid commits, premise
acceptance, lineage loss --- that are invisible to downstream answer
accuracy?

\textbf{RQ3 --- Measurement.} Which metrics separate decision
correctness, memory-state correctness, and downstream answer behavior,
and which shipped metrics fail (and must be disclosed as failing)?

\textbf{RQ4 --- Reliability.} How far do the reported conclusions
survive paired significance testing with multiplicity control,
annotation audit, and adversarial review of run provenance?

\textbf{RQ5 --- Reproducibility.} What artifact and claim controls keep
a thesis, its papers, and its code mutually consistent --- including
when the honest answer is ``this number is not final''?

\section{1.4 Hypotheses and their outcomes
(preview)}\label{hypotheses-and-their-outcomes-preview}

The original directional comparisons are not confirmatory because the
pilot fails its shortcut gate. Chapter 7 therefore treats them as
diagnostics and reports these outcomes:

\begin{itemize}
\tightlist
\item
  \textbf{H1 --- Governance framing helps an LLM.} Not identified by v2:
  structured and unstructured prompts differ on a label-leaking
  collection, so the observed gap is descriptive.
\item
  \textbf{H2 --- An interpretable reference policy can match an LLM
  policy.} No superiority or non-inferiority conclusion: test accuracy
  is 0.8889 versus 0.9000 and the family-clustered paired interval
  includes zero.
\item
  \textbf{H3 --- Abstention is load-bearing for action safety.}
  Mechanistically supported inside the generator, but not validated
  externally; zero observed action false commits has a non-zero
  finite-sample upper bound and accompanies missed updates.
\item
  \textbf{H4 --- Decision and stored-state metrics can diverge.}
  Supported as a measurement demonstration: the no-lineage ablation
  preserves action accuracy while degrading lineage integrity.
\item
  \textbf{H5 --- Controller end-task superiority on LongMemEval.} Not
  supported; existing downstream runs are contaminated, incomplete, or
  selection-biased.
\end{itemize}

\section{1.5 Contributions}\label{contributions}

\begin{enumerate}
\def\labelenumi{\arabic{enumi}.}
\tightlist
\item
  \textbf{Formulation and evaluation decomposition (RQ1, RQ3).} Memory
  writes are represented as KEEP/SUPERSEDE/ABSTAIN decisions over typed,
  provenance-tracked, lineage-preserving records and evaluated
  separately at action, state, and answer layers.
\item
  \textbf{A benchmark with build-time construct gates (RQ2, RQ4).} MGP-6
  v4.3 --- 900 human-authored items in 100 semantic components across 20
  domains, component-split, exactly action-balanced, dual-annotated,
  with shortcut/duplicate/allowlist/deterministic-build gates enforced
  at release rather than audited afterwards, and a sealed 540-item
  hidden partition governed by a one-shot pre-registration protocol.
\item
  \textbf{A confirmatory decision-level result inside stated bounds
  (RQ1).} One pre-registered configuration scored once on the sealed
  split: 0.8481 accuracy, 0.8479 macro-F1, Cohen's κ 0.7722, with
  component-cluster intervals roughly 40\% wider than item-independence
  intervals would imply.
\item
  \textbf{Error decomposition by operational harm (RQ3).} The thesis
  separates state-preserving from state-corrupting write failures and
  shows the distinction changes the risk estimate fivefold (0.2111
  aggregate false commit versus 0.0444 state-corrupting), alongside
  missed-update, destructive-overwrite, and unnecessary-abstention rates
  each with finite-sample intervals.
\item
  \textbf{Component attribution (RQ2).} Scoring the rule policy and the
  model policy separately shows neither exceeds 0.750 alone against a
  composed 0.8444, and bounds the abstention guard's hidden-split
  contribution at most 6.7 accuracy points of the 84.8 observed.
\item
  \textbf{Computation and cost accounting (RQ3).} Per-decision latency
  distributions, prompt/completion token counts, retry and failure
  counts, and cost-conditional expected loss under three governance
  regimes --- with monetary cost declared unavailable rather than
  reconstructed from list prices.
\item
  \textbf{Measured negative results (RQ4).} The v2 pilot's conflict-type
  shortcut reaches 1.000 held-out accuracy; the v2 reference policy does
  not outperform the structured LLM and its preferred status flips with
  the cost matrix; and the v4.3 controller's abstentions are not
  selectively risk-targeted. Each is reported rather than omitted.
\item
  \textbf{Downstream validation with an attributable pattern (RQ1,
  RQ5).} A homogeneous full-slice LongMemEval-M run places governed
  writes at 76.1\% above every structural baseline, with every adjacent
  step in the ordering separating under paired tests and the gain
  concentrated in the competencies the write policy targets.
\item
  \textbf{Evidence exclusion and successor design (RQ4, RQ5).}
  Incomplete, mapped, contaminated, and selection-biased results are
  excluded from primary claims by an explicit ledger; a role-separated
  successor protocol specifies independent generation, three-human
  annotation, long trajectories, hidden baselines scored under the same
  one-shot discipline, and native external metrics.
\item
  \textbf{Reproducible evidence package (RQ5).} Released per-item
  predictions, a recomputation artifact regenerating every reported
  number from raw predictions, audit JSON/Markdown, claim-evidence
  matrix, and machine-readable ledgers keep paper, thesis, code, and
  limitations synchronized.
\end{enumerate}

\section{1.6 Scope}\label{scope}

The thesis evaluates the ASK formulation on an author-designed
diagnostic slice and one stratified LongMemEval-M slice with a single
base model family and single seed. It does not claim community-benchmark
status for MGP-6 v2, does not claim end-task superiority on LongMemEval,
does not claim human multi-annotator agreement, and does not claim that
governance on these slices generalizes to open-domain deployments. These
boundaries are stated because they are what the evidence supports;
Chapter 9 enumerates them systematically.

\section{1.7 Thesis structure and contribution
map}\label{thesis-structure-and-contribution-map}

{\def\LTcaptype{none} % do not increment counter
\begin{longtable}[]{@{}
  >{\raggedright\arraybackslash}p{(\linewidth - 4\tabcolsep) * \real{0.3333}}
  >{\raggedright\arraybackslash}p{(\linewidth - 4\tabcolsep) * \real{0.3333}}
  >{\raggedright\arraybackslash}p{(\linewidth - 4\tabcolsep) * \real{0.3333}}@{}}
\toprule\noalign{}
\begin{minipage}[b]{\linewidth}\raggedright
Research question
\end{minipage} & \begin{minipage}[b]{\linewidth}\raggedright
Main chapter(s)
\end{minipage} & \begin{minipage}[b]{\linewidth}\raggedright
Primary artifact
\end{minipage} \\
\midrule\noalign{}
\endhead
\bottomrule\noalign{}
\endlastfoot
RQ1 & Chapter 4 & \texttt{src/ask/} (policy, memory operations,
orchestrator) \\
RQ2 & Chapters 5, 7 & \texttt{benchmarks/mgp6/} (data, manual, metrics),
ablation runs \\
RQ3 & Chapters 6, 7 & \texttt{metrics\_mgp6.py}, ledger
\texttt{known\_metric\_issues} \\
RQ4 & Chapters 7, 8, 9 & \texttt{reports/STATISTICAL\_ANALYSIS.md},
\texttt{reports/ERROR\_TAXONOMY.md},
\texttt{reports/THREATS\_TO\_VALIDITY.md} \\
RQ5 & Chapters 6, 9, App. E--F &
\texttt{artifacts/results\_ledger.json},
\texttt{reports/CLAIMS\_EVIDENCE\_MATRIX.md} \\
\end{longtable}
}
